% !TEX root = ../main.tex
% Abstrak (pedoman 10.17, format monograf):
%   - satu paragraf, maksimal 500 kata, satu spasi;
%   - dua bahasa, sebisa mungkin dalam satu halaman;
%   - nama (tanpa NPM), program studi, dan judul di atas abstrak;
%   - kata kunci di bawah, dalam bahasa masing-masing;
%   - semua istilah asing (kecuali nama) dicetak miring.
\chapter*{ABSTRAK}
\phantomsection\addcontentsline{toc}{chapter}{ABSTRAK}
\vspace{-1\baselineskip}
{\setstretch{1}
\noindent\begin{tabular}{@{}p{3cm}@{\,:\,}>{\raggedright\arraybackslash}p{10.5cm}@{}}
Nama & \thenama \\
Program Studi & \theprodi \\
Judul & \thejudul \\
Pembimbing & \thepembimbing \\
\end{tabular}

\vspace{0.4cm}
\noindent Tulis abstrak dalam satu paragraf: masalah dan celah penelitian, tujuan,
metode (konfigurasi \textit{single-model}, \textit{self-consistency}, dan
\textit{LLM council} pada Agen Verifikasi), data, hasil utama dengan angka, dan
kesimpulan. Maksimal 500 kata.

\vspace{0.3cm}
\noindent\textbf{Kata kunci:} kata kunci 1, kata kunci 2, kata kunci 3

\vspace{1cm}
\begin{center}\bfseries ABSTRACT\end{center}
\noindent\begin{tabular}{@{}p{3cm}@{\,:\,}>{\raggedright\arraybackslash}p{10.5cm}@{}}
Name & \thenama \\
Study Program & \theprodi \\
Title & \thejudulinggris \\
Supervisor & \thepembimbing \\
\end{tabular}

\vspace{0.4cm}
\noindent\textit{Write the English abstract here, as a single paragraph of at most
500 words, matching the Indonesian abstract.}

\vspace{0.3cm}
\noindent\textbf{\textit{Keywords:}} \textit{keyword 1, keyword 2, keyword 3}
\par}
